\section{Preliminaries}

\subsection{Properties of the Bargmann-Fock space} We briefly recall the connection between the Gaussian short-time Fourier transform
(STFT) and the Bargmann--Fock space, which will be used systematically in what
follows. Excellent references for the material in this subsection are
\cite{Bargmann,Folland,Grochenig,Zhu}; see also
\cite{Berezin,Berger-Coburn,Daubechies}.

\medskip

Let $\varphi(x) = 2^{1/4} e^{-\pi x^2}$ be the normalized Gaussian. For $f \in L^2(\mathbb{R})$, its Gaussian STFT is
\begin{equation}\label{eq:STFT}
V_\varphi f(x,\omega)
= \int_{\mathbb{R}} f(t)\, \varphi(x-t)\, e^{-2\pi i t \omega}\, dt.
\end{equation}
The Bargmann transform
$\mathcal{B} : L^2(\mathbb{R}) \to \mathcal{F}^2(\mathbb{C})$ is defined by
\begin{equation}\label{eq:Bargmann}
(\mathcal{B}f)(z)
= 2^{1/4} \int_{\mathbb{R}} f(t)\, e^{2\pi t z - \pi t^2 - \frac{\pi}{2} z^2}\, dt,
\qquad z \in \mathbb{C}.
\end{equation}
It is well known that $\mathcal{B}$ is a unitary isomorphism onto the Bargmann--Fock
space $\mathcal{F}^2(\mathbb{C})$, endowed with the norm
\[
\|F\|_{\mathcal{F}^2}^2
= \int_{\mathbb{C}} |F(z)|^2 e^{-\pi |z|^2} \, dA(z).
\]
A direct computation shows that, for $z = x + i\omega$,
\begin{equation}\label{eq:STFT-B}
V_\varphi f(x,-\omega)
= e^{\pi i x \omega}\, (\mathcal{B}f)(z)\, e^{-\frac{\pi}{2}|z|^2}.
\end{equation}
In particular, setting $F = \mathcal{B}f$, one obtains
\begin{equation}\label{eq:energy}
\int_{\mathbb{R}^2} |V_\varphi f(x,\omega)|^2 \, dx\, d\omega
= \int_{\mathbb{C}} |F(z)|^2 e^{-\pi |z|^2} \, dA(z)
= \|F\|_{\mathcal{F}^2}^2.
\end{equation}

\medskip

\noindent
\subsubsection*{Weyl operators and invariance.}
For $a \in \mathbb{C}$, define the operator
\begin{equation}\label{eq:Ta}
(T_a F)(z) := e^{\pi \overline{a} z - \frac{\pi}{2}|a|^2} F(z - a).
\end{equation}
Then $T_a$ is unitary on $\mathcal{F}^2(\mathbb{C})$. Indeed,
\begin{align*}
\|T_aF\|_{\mathcal F^2}^2
&=
\int_{\C} e^{2\pi \Re(\overline{a}z)-\pi |a|^2}\,|F(z-a)|^2 e^{-\pi |z|^2}\,dA(z) \\
&=
\int_{\C} |F(z-a)|^2 e^{-\pi |z-a|^2}\,dA(z)
=
\|F\|_{\mathcal F^2}^2,
\end{align*}
because
\[
2\Re(\overline{a}z)-|a|^2-|z|^2=-|z-a|^2.
\]
Moreover,
\begin{equation}\label{eq:uF-translate}
u_{T_a F}(z) = u_F(z - a).
\end{equation}
Indeed, since
\(2\Re(\overline{a}z)-|a|^2-|z|^2 = -|z-a|^2\),
\begin{equation}\label{eq:u-translation-global}
u_{T_aF}(z)
=
|T_aF(z)|^2e^{-\pi |z|^2}
=
e^{2\pi\Re(\overline{a}z)-\pi|a|^2}|F(z-a)|^2e^{-\pi |z|^2}
=
|F(z-a)|^2e^{-\pi |z-a|^2}=u_F(z-a).
\end{equation}
In particular, the functionals $I_F$ defined by,
for \(s>0\),
\[
I_F(s):=\sup_{|\Omega|=s}\int_\Omega u_F\,dA,
\]
satisfy, by the translation-invariance of the Lebesgue measure on \(\C\),
\begin{equation}\label{eq:I-translation-global}
I_{T_a F}(s) = I_F(s), \qquad a \in \mathbb{C},
\end{equation}
reflecting the invariance of the Gaussian STFT under time--frequency shifts;
see \cite{Daubechies,Grochenig,Kulikov-Oliveira-Ramos,Goncalves-Oliveira-Steinerberger}.

\medskip

\noindent
\subsubsection*{Connection with localization.}
If $F = \mathcal{B}f$, then for every measurable $\Omega \subset \mathbb{C}$,
\[
\int_\Omega u_F \, dA
= \int_{\widetilde{\Omega}} |V_\varphi f(x,\omega)|^2 \, dx\, d\omega,
\]
where $\widetilde{\Omega} = \{(x,\omega) : (x,-\omega) \in \Omega\}$.
Thus, concentration properties of the STFT are naturally encoded by the function
\begin{equation}\label{eq:uF}
u_F(z) := |F(z)|^2 e^{-\pi |z|^2}, \qquad z \in \mathbb{C},
\end{equation}
which is the main object of interest in the Bargmann--Fock setting, in terms of the
quantities $I_F(s)$.

\medskip

We next collect a few elementary facts that will be used repeatedly.

\begin{lemma}\label{lem:Fock-basic}
Let $F,G \in \mathcal{F}^2(\mathbb{C})$. Then:
\begin{enumerate}
\item[(i)] The space $\mathcal{F}^2(\mathbb{C})$ is a reproducing kernel Hilbert space with kernel
\[
K_w(z) = e^{\pi z \overline{w}}, \qquad z,w \in \mathbb{C}.
\]

\item[(ii)] (Sharp pointwise bound) For every $z \in \mathbb{C}$,
\[
u_F(z) \le \|F\|_{\mathcal{F}^2}^2,
\]
and equality holds at some point $w \in \mathbb{C}$ if and only if $F$ is
proportional to the reproducing kernel $K_w$.

\item[(iii)] $u_F(z) \to 0$ as $|z| \to \infty$.

\item[(iv)] One has the estimate
\[
\|u_F - u_G\|_{L^1(\mathbb{C})}
\le \big( \|F\|_{\mathcal{F}^2} + \|G\|_{\mathcal{F}^2} \big)
\|F - G\|_{\mathcal{F}^2}.
\]
\end{enumerate}
\end{lemma}

\begin{proof}
(i) For \(w\in \C\), the function \(K_w(z)=e^{\pi z\overline w}\) is entire and
satisfies, after completing the square via
\(2\Re(z\overline w)-|z|^2=-|z-w|^2+|w|^2\),
\[
\|K_w\|_{\mathcal F^2}^2
=
\int_{\C} e^{2\pi \Re(z\overline w)-\pi |z|^2}\,dA(z)
=
e^{\pi |w|^2}\int_{\C}e^{-\pi |z-w|^2}\,dA(z)
=
e^{\pi |w|^2},
\]
since the Gaussian \(e^{-\pi |\cdot|^2}\) integrates to \(1\) over \(\C\) (with
the Lebesgue measure \(dA\) on \(\C\simeq \mathbb R^2\)). In particular
\(K_w\in \mathcal F^2(\C)\). Now, expanding \(F\in \mathcal F^2(\C)\) along the
orthonormal basis \(\{e_k\}_{k\geq 0}\) of Lemma \ref{lem:monomials} and using
the Taylor expansion
\(K_w(z)=\sum_{k\geq 0}\frac{\pi^k\overline w^k}{k!}z^k=\sum_{k\geq 0}\overline{e_k(w)}\,e_k(z)\),
one obtains the reproducing identity
\[
\langle F,K_w\rangle_{\mathcal F^2}
=
\sum_{k\geq 0} a_k\,e_k(w)
=
F(w),
\qquad F(z)=\sum_{k\geq 0}a_k e_k(z).
\]
Point evaluation is therefore continuous, and \(F\) is entire as a uniform
limit on compact sets of polynomials. This is the standard description of the
Bargmann--Fock space, see \cite{Bargmann,Folland,Grochenig,Zhu}.

(ii) Applying Cauchy--Schwarz to the reproducing identity from (i),
\begin{equation}\label{eq:pointwise}
|F(z)|
=
|\langle F,K_z\rangle_{\mathcal F^2}|
\le
\|F\|_{\mathcal F^2}\,\|K_z\|_{\mathcal F^2}
=
\|F\|_{\mathcal F^2}\,e^{\frac{\pi}{2}|z|^2},
\end{equation}
which immediately gives \(u_F(z)=|F(z)|^2 e^{-\pi |z|^2}\le \|F\|_{\mathcal F^2}^2\).
Equality at a point \(z=w\) corresponds to equality in Cauchy--Schwarz, that is,
\(F=c\,K_w\) for some \(c\in \C\); conversely, for \(F=c K_w\) one computes
\(u_F(w)=|c|^2 e^{2\pi |w|^2}e^{-\pi |w|^2}=|c|^2 e^{\pi |w|^2}=|c|^2\|K_w\|_{\mathcal F^2}^2
=\|F\|_{\mathcal F^2}^2\), proving the rigidity statement. The sharpness of this
pointwise bound goes back to \cite{Lieb}.

(iii) Fix \(\varepsilon>0\), and choose a polynomial \(P\) such that
\[
\|F-P\|_{\mathcal F^2}<\varepsilon
\]
using the density of polynomials in \(\mathcal F^2(\C)\). By the pointwise bound from (ii),
\[
|F(z)-P(z)|e^{-\frac{\pi}{2}|z|^2}\le \varepsilon
\qquad\text{for every }z\in \C.
\]
Since \(P\) is a polynomial,
\[
P(z)e^{-\frac{\pi}{2}|z|^2}\to 0
\qquad\text{as }|z|\to\infty.
\]
Hence
\[
\limsup_{|z|\to\infty}|F(z)|e^{-\frac{\pi}{2}|z|^2}\le \varepsilon.
\]
As \(\varepsilon>0\) is arbitrary, we conclude that
\[
|F(z)|e^{-\frac{\pi}{2}|z|^2}\to 0,
\]
that is, \(u_F(z)\to 0\) as \(|z|\to\infty\).

(iv) Writing
\[
|u_F - u_G|
= \big||F|^2 - |G|^2\big| e^{-\pi|z|^2}
\le |F-G|(|F|+|G|) e^{-\pi|z|^2},
\]
the claim follows by Cauchy--Schwarz.
\end{proof}

\medskip

We now note the following crucial fact about the Bargmann transform: it
transforms the basis of normalized Hermite functions \(\{h_k\}_{k\geq 0}\) on
\(L^2(\R)\) (see \cite{Thangavelu}) onto normalized monomials, yielding a
natural orthonormal basis on $\mathcal{F}^2(\C).$ More precisely, one has
\(\mathcal B h_k=e_k\) for every \(k\ge 0\), see \cite{Bargmann,Folland,Grochenig}.

\begin{lemma}\label{lem:monomials}
The functions
\[
e_k(z) := b_k z^k, \qquad k \ge 0,
\]
form an orthonormal basis of $\mathcal{F}^2(\mathbb{C})$, where
\[
b_k = \left( \frac{\pi^k}{k!} \right)^{1/2}.
\]
In particular,
\[
\langle e_k, e_m \rangle_{\mathcal{F}^2} = \delta_{k,m}.
\]
\end{lemma}

\begin{proof}
Writing \(z=re^{i\theta}\) and using \(dA(z)=r\,dr\,d\theta\) on \(\C\),
a direct computation in polar coordinates gives
\[
\int_{\mathbb{C}} |z|^{2k} e^{-\pi |z|^2}\,dA(z)
= \int_0^{2\pi}\!\!\int_0^\infty r^{2k+1} e^{-\pi r^2}\,dr\,d\theta
= 2\pi \int_0^\infty r^{2k+1} e^{-\pi r^2}\,dr
= \frac{k!}{\pi^k},
\]
where the last identity follows from the substitution \(s=\pi r^2\), giving
\(2\pi\int_0^\infty r^{2k+1}e^{-\pi r^2}\,dr=\pi^{-k}\int_0^\infty s^k e^{-s}\,ds
=\pi^{-k}\Gamma(k+1)\). Thus
\(\|z^k\|_{\mathcal{F}^2}^2 = \frac{k!}{\pi^k}\), which yields the normalization.
For orthogonality, again writing \(z=re^{i\theta}\),
\[
\langle z^k,z^m\rangle_{\mathcal F^2}
=
\int_0^\infty r^{k+m+1}e^{-\pi r^2}\,dr\int_0^{2\pi}e^{i(k-m)\theta}\,d\theta
=
0
\qquad\text{for }k\neq m.
\]
For completeness, since \(\mathcal B\colon L^2(\R)\to \mathcal F^2(\C)\) is a
unitary isomorphism and \(\mathcal Bh_k=e_k\) for every \(k\ge 0\), where
\(\{h_k\}_{k\ge 0}\) is the orthonormal basis of Hermite functions on
\(L^2(\R)\) (see \cite{Thangavelu}), the family \(\{e_k\}_{k\ge 0}\) is the
image of an orthonormal basis under a unitary, hence itself an orthonormal
basis of \(\mathcal F^2(\C)\). Equivalently, polynomials are dense in
\(\mathcal F^2(\C)\); see \cite{Bargmann,Zhu}.
\end{proof}

\medskip

\begin{lemma}\label{lem:growth}
Let $F \in \mathcal{F}^2(\mathbb{C})$. Then:
\begin{enumerate}
\item[(i)] $F$ admits the expansion
\[
F(z) = \sum_{k=0}^\infty a_k e_k(z),
\qquad \sum_{k=0}^\infty |a_k|^2 = \|F\|_{\mathcal{F}^2}^2.
\]

\item[(ii)] (Growth bound) For every $z \in \mathbb{C}$,
\[
|F(z)| \le \|F\|_{\mathcal{F}^2} e^{\frac{\pi}{2}|z|^2}.
\]

\item[(iii)] (Derivative bounds) For every $z \in \mathbb{C}$,
\[
|F'(z)| \le C(1+|z|)\, \|F\|_{\mathcal{F}^2} e^{\frac{\pi}{2}|z|^2}.
\]
\end{enumerate}
\end{lemma}

\begin{proof}
(i) follows from the orthonormal basis in Lemma \ref{lem:monomials}.

(ii) follows from the reproducing kernel estimate:
\[
|F(z)| = |\langle F, K_z \rangle|
\le \|F\|_{\mathcal{F}^2} \|K_z\|_{\mathcal{F}^2}
= \|F\|_{\mathcal{F}^2} e^{\frac{\pi}{2}|z|^2}.
\]

(iii) Starting from the reproducing identity
\(F(z)=\langle F,K_z\rangle_{\mathcal F^2}=\int_{\C}F(w)e^{\pi \overline{w}z}
e^{-\pi |w|^2}\,dA(w)\) and differentiating in \(z\) under the integral sign
(justified by the absolute convergence on compacta of \(z\), via the pointwise
bound from (ii)), we obtain
\[
F'(z)=\pi \int_{\C}F(w)\,\overline{w}\,e^{\pi\overline{w}z}e^{-\pi |w|^2}\,dA(w)
=\pi \langle F, w\,K_z\rangle_{\mathcal F^2},
\]
since \(\overline{wK_z(w)}=\overline{w}\,e^{\pi\overline{w}z}\). Cauchy--Schwarz
then gives
\[
|F'(z)|
\le
\pi \|F\|_{\mathcal F^2}\,\|wK_z\|_{\mathcal F^2}.
\]
Now, by completing the square as in (i),
\[
\|wK_z\|_{\mathcal F^2}^2
=
\int_{\C} |w|^2 e^{2\pi\Re(w\overline z)-\pi |w|^2}\,dA(w)
=
e^{\pi |z|^2}\int_{\C} |w+z|^2 e^{-\pi |w|^2}\,dA(w),
\]
after the change of variables \(w\mapsto w+z\). Expanding
\(|w+z|^2=|w|^2+2\Re(w\overline z)+|z|^2\), and using that
\(\int_\C e^{-\pi |w|^2}\,dA(w)=1\),
\(\int_\C |w|^2 e^{-\pi |w|^2}\,dA(w)=\frac{1}{\pi}\), and
\(\int_\C w\,e^{-\pi |w|^2}\,dA(w)=0\) (by rotational invariance), we conclude
\[
\int_{\C} |w+z|^2 e^{-\pi |w|^2}\,dA(w)
=
\frac{1}{\pi}+|z|^2
\le
C(1+|z|^2),
\]
for some absolute constant \(C>0\). Therefore
\[
|F'(z)| \le C(1+|z|)\,\|F\|_{\mathcal F^2}e^{\frac{\pi}{2}|z|^2}.
\]
\end{proof}

\medskip

\begin{lemma}\label{lem:max}
Let $F \in \mathcal{F}^2(\mathbb{C})$ and set
\[
T := \sup_{z \in \mathbb{C}} u_F(z).
\]
Then $0 \le T \le \|F\|_{\mathcal{F}^2}^2$, and $T = \|F\|_{\mathcal{F}^2}^2$
is attained if and only if $F$ is proportional to the reproducing kernel $K_w$
for some $w \in \mathbb{C}$.
\end{lemma}

\begin{proof}
The bound and the characterization of the equality case are immediate from
Lemma~\ref{lem:Fock-basic}(ii).
    \end{proof}

\medskip

    We proceed with a few elementary facts.

\begin{lemma}\label{lem:basic-properties-global}
Let \(F,G\in \mathcal F^2(\C)\). Then:
\begin{enumerate}
    \item[(i)] \(I_F(s)\) is attained by a superlevel set of \(u_F\), and
    \[
    0\le I_F(s)\le \|F\|_{\mathcal F^2}^2.
    \]
    \item[(ii)] For every \(c\in \C\),
    \[
    I_{cF}(s)=|c|^2 I_F(s).
    \]
    \item[(iii)] One has the Lipschitz bound
    \[
    |I_F(s)-I_G(s)|
    \le
    \|u_F-u_G\|_{L^1(\C)}
    \le
    \big(\|F\|_{\mathcal F^2}+\|G\|_{\mathcal F^2}\big)\,
    \|F-G\|_{\mathcal F^2}.
    \]
\end{enumerate}
\end{lemma}

\begin{proof}
The identity
\[
\int_\C u_F\,dA=\int_\C |F(z)|^2e^{-\pi |z|^2}\,dA(z)=\|F\|_{\mathcal F^2}^2
\]
is immediate from the definition of the Fock norm. By
Lemma~\ref{lem:Fock-basic}(ii)--(iii), \(u_F\) is bounded, continuous, and
satisfies \(u_F(z)\to 0\) as \(|z|\to\infty\); thus
\(u_F\in C_0(\C)\cap L^1(\C)\).

The existence of a maximizing set for \(I_F(s)\), of the form
\(\{u_F>t\}\) (possibly together with part of \(\{u_F=t\}\)) with \(t\geq 0\)
chosen so that the resulting set has measure \(s\), is the classical bathtub
principle, which applies since \(u_F\in C_0(\C)\cap L^1(\C)\). The bound
\[
0\le I_F(s)\le \int_\C u_F\,dA=\|F\|_{\mathcal F^2}^2
\]
is immediate. The scaling identity in (ii) is obvious.

For (iii), for every measurable \(\Omega\subset \C\) with \(|\Omega|=s\),
\[
\left|\int_\Omega u_F\,dA-\int_\Omega u_G\,dA\right|
\le \int_\C |u_F-u_G|\,dA.
\]
Taking suprema over \(|\Omega|=s\) gives
\[
|I_F(s)-I_G(s)|\le \|u_F-u_G\|_{L^1(\C)}.
\]
Finally,
\begin{align*}
\|u_F-u_G\|_{L^1(\C)}
&=
\int_\C \big||F|^2-|G|^2\big|\,e^{-\pi |z|^2}\,dA \\
&\le
\int_\C |F-G|(|F|+|G|)e^{-\pi |z|^2}\,dA \\
&\le
\|F-G\|_{\mathcal F^2}\big(\|F\|_{\mathcal F^2}+\|G\|_{\mathcal F^2}\big)
\end{align*}
by Cauchy--Schwarz.
\end{proof}

